\documentclass[11pt]{amsart}
\usepackage{amsmath,amssymb,amsthm,booktabs}
\usepackage[margin=1.15in]{geometry}
\usepackage{hyperref}
\newtheorem{theorem}{Theorem}
\newtheorem{proposition}[theorem]{Proposition}
\newtheorem{lemma}[theorem]{Lemma}
\newtheorem{corollary}[theorem]{Corollary}
\theoremstyle{definition}
\newtheorem{remark}[theorem]{Remark}
\newcommand{\Prob}{\mathbb{P}}
\newcommand{\E}{\mathbb{E}}
\newcommand{\tmix}{\tau_1^{(2)}}
\newcommand{\hE}{\widehat{\E}}
\newcommand{\hw}{\widehat{w}}
\newcommand{\hpi}{\widehat{\pi}}
\newcommand{\hI}{\widehat{I}}
\newcommand{\hA}{\widehat{A}}
\newcommand{\hX}{\widehat{X}}
\newcommand{\lra}{\leftrightarrow}

\title[Lumping and the mixing time of a reversible chain]{Lumping and the mixing time of a reversible chain:
the dependence on the number of states}
\author{Yitzchak Shmalo}
\date{November 2026}

\begin{document}
\begin{abstract}
Aldous and Fill asked whether the stopping-rule mixing time $\tmix$ of a chain obtained from a reversible Markov
chain by lumping states is at most an absolute constant times that of the original chain. Hermon and Peres answered
the question in the negative: for random walks on graphs of bounded degree, lumping can increase the mixing time by a
factor of the order of the logarithm of the number of vertices, and by no more. We consider general reversible chains
with $n$ states and show that the largest possible factor $K(n)$ satisfies
$\lfloor (n-2)/2\rfloor\le K(n)\le n-1$. The upper bound follows from standard inequalities and from the inequality
$\tau_2\le\tmix$ between the relaxation time and the stopping-rule mixing time, for which we give a short proof;
the lower bound comes from an explicit random walk on a weighted comb. We also show that lumping increases the mean
hitting time of a set $A$ by a factor of at most $n-|A|$, and that this bound cannot be improved by more than a factor
of about two.
\end{abstract}
\maketitle

\section{Introduction}
Let $I$ be a finite set with $n\ge2$ elements and let $(w_{xy})_{x,y\in I}$ be a symmetric array of non-negative
weights, loops $w_{xx}>0$ being allowed, such that the graph with edges $\{x,y\}$, $w_{xy}>0$, is connected. Put
$w_x=\sum_y w_{xy}$ and $w=\sum_x w_x$. The random walk on this network is the Markov chain $X$ on $I$ with transition
probabilities $p_{xy}=w_{xy}/w_x$; it is reversible with stationary distribution $\pi_x=w_x/w$, and every irreducible
reversible chain on $I$ is of this form. We write $T_y=\inf\{t\ge0:X_t=y\}$ for the hitting time of a state $y$,
$T_A=\inf\{t\ge0:X_t\in A\}$ for the hitting time of a set $A$, so that $T_y=0$ for the walk started at $y$, and $\E_\mu$ for expectation for the chain started from the distribution $\mu$.

Given a map $f$ from $I$ onto a set $\hI$, the \emph{induced chain} $\hX$ is the random walk on $\hI$ with weights
\[
\hw_{ab}=\sum_{x\in f^{-1}(a)}\ \sum_{y\in f^{-1}(b)} w_{xy},\qquad a,b\in\hI .
\]
Its transition probabilities are $\Prob_\pi(f(X_1)=b\mid f(X_0)=a)$ and its stationary distribution is
$\hpi_a=\pi(f^{-1}(a))$; this is the chain of Aldous and Fill~\cite[Section~4.6.1]{AF}. In the language of networks,
the vertices with the same value of $f$ are glued together. We say that $\hX$ is obtained from $X$ by lumping, and
write $\hE$ and $\widehat{T}$ for its expectations and hitting times.

The parameter of~\cite[Chapter~4]{AF} that we study is
\[
\tmix=\max_{i\in I}\ \min_{U}\ \E_iU ,
\]
where the minimum is over all stopping times $U$, possibly randomized with a randomization independent of the future
of the chain, such that $X_U$ has distribution $\pi$ under
$\Prob_i$. Lov\'asz and Winkler~\cite{LW95} showed that the minimum is attained and found its value
(see also~\cite[Theorem~9.4]{AF}), so that
\begin{equation}\label{eq:lw}
\tmix=\max_{i,j\in I}\bigl(\E_iT_j-\E_\pi T_j\bigr).
\end{equation}
For reversible chains $\tmix$ is equivalent, up to absolute constants, to the total variation mixing time of the
chain in continuous time~\cite{Aldous82}, \cite[Theorem~4.6]{AF}; see also~\cite{LW98}.

The maximal mean commute time $\tau^*=\max_{i,j}(\E_iT_j+\E_jT_i)$, the average hitting time
$\tau_0=\sum_j\pi_j\E_iT_j$ (which does not depend on $i$), the relaxation time $\tau_2=1/(1-\lambda_2)$, where
$\lambda_2$ is the second largest eigenvalue of $(p_{xy})$, and the Cheeger time constant $\tau_c$ do not increase
under lumping~\cite[Proposition~4.44]{AF}. For the mixing time the matter was left open:

\smallskip
\noindent\textbf{Open Problem 4.45 of~\cite{AF}.} \emph{Is the value of $\tmix$ in an induced chain bounded by $K$
times the value of $\tmix$ in the original chain, for some absolute constant $K$? For $K=1$?}
\smallskip

Hermon and Peres~\cite[Theorem~3(c)]{HP} answered both questions in the negative. They construct graphs $G_n$ of
bounded degree for which lumping together certain pairs of vertices increases the total variation mixing time of the
lazy simple random walk by a factor of order $\log|V(G_n)|$. Lumping does not increase the relaxation time, and the
mixing time of a lazy reversible chain is bounded below by a constant multiple of the relaxation time and above by a
constant multiple of the relaxation time multiplied by $\log(1/\min_x\pi_x)$; hence this is the largest possible
order in that setting~\cite[Remark~1.8]{HP}.

For a random walk on a weighted network $\min_x\pi_x$ can be arbitrarily small, and the bound in terms of
$\log(1/\min_x\pi_x)$ says nothing about the dependence on the number of states. This note determines that dependence
up to a factor of about two. Let
\[
K(n)=\sup\ \frac{\tmix(\hX)}{\tmix(X)} ,
\]
the supremum being over all reversible chains $X$ with at least two and at most $n$ states and all induced chains
$\hX$ of $X$.

\begin{theorem}\label{thm:main}
For every $n\ge2$,
\[
\max\Bigl(1,\Bigl\lfloor\frac{n-2}{2}\Bigr\rfloor\Bigr)\ \le\ K(n)\ \le\ n-1 .
\]
\end{theorem}

The upper bound is a chain of four inequalities, three of which are in~\cite{AF}. The fourth is the following
comparison of the relaxation time with the stopping-rule mixing time, in which the constant is one; whether this holds
is asked in~\cite[Open Problem~4.48(b)]{AF}. Comparisons of the two quantities up to absolute constants are
in~\cite{Aldous82,LW98}.

\begin{lemma}\label{lem:relax}
Let $\lambda<1$ be a real eigenvalue of the transition matrix of an irreducible chain on a finite set, with a real
eigenfunction $g$, and let $i$ be a state at which $|g|$ is largest. If $U$ is a stopping time such that $X_U$ has the
stationary distribution $\pi$ under $\Prob_i$, then $\E_iU\ge 1/(1-\lambda)$. In particular $\tau_2\le\tmix$ for
every reversible chain.
\end{lemma}

Equality $\tau_2=\tmix$ holds for every chain with two states and for the simple random walk on a complete graph.

The lower bound in Theorem~\ref{thm:main} is approached along a family of random walks on a weighted comb with $2N+2$
states, described in Section~\ref{sec:comb}; lumping each tooth of the comb with a vertex of the spine multiplies
$\tmix$ by a factor that can be made arbitrarily close to $N$. In these examples the stationary distribution is very
far from uniform, in agreement with the bound of~\cite{HP}; see Remark~\ref{rem:pimin}.

The mechanism is a one-sided hitting time. Aldous and Fill observe~\cite[Section~4.6.1]{AF} that, although commute
times decrease under lumping, a mean hitting time may increase, and give an example on four vertices where it
doubles. The next proposition bounds the increase for hitting times of sets.

\begin{proposition}\label{prop:hit}
Let $\hX$ be the chain induced from $X$ by $f:I\to\hI$, let $\hA$ be a non-empty proper subset of $\hI$ and
$A=f^{-1}(\hA)$. Then for every $a\in\hI\setminus\hA$
\[
\hE_a\widehat{T}_{\hA}\ \le\ \sum_{y\in I\setminus A}\E_yT_A\ \le\ (n-|A|)\max_{y\in I}\E_yT_A .
\]
For every $N\ge1$ and $\varepsilon>0$ there is a chain with $n=2N+2$ states, an induced chain and a set $A$
with one element such that $\hE_a\widehat{T}_{\hA}\ge(N-\varepsilon)\max_y\E_yT_A$ for some $a$.
\end{proposition}

Lemma~\ref{lem:relax} and the upper bound of Theorem~\ref{thm:main} are proved in Section~\ref{sec:upper}, the lower
bound in Section~\ref{sec:comb} and Proposition~\ref{prop:hit} in Section~\ref{sec:hit}. We do not know whether
$K(n)/n$ converges, nor the value of $K(4)$; an example with four states and ratio above $1.27$ is in
Remark~\ref{rem:four}.

\section{The upper bound}\label{sec:upper}

\begin{proof}[Proof of Lemma~\ref{lem:relax}]
Replacing $g$ by a multiple we may assume that $g(i)=1=\max_x|g(x)|$. Since $Pg=\lambda g$, the process
\[
M_t=g(X_t)+(1-\lambda)\sum_{s<t}g(X_s),\qquad t\ge0,
\]
is a martingale with $M_0=1$ under $\Prob_i$, with respect to the filtration generated by the chain and by the
randomization used by $U$. We may assume that $\E_iU<\infty$. The increments of $M$ are bounded by $2$, so the optional
stopping theorem gives $\E_iM_U=1$. As $\lambda\ne1$, the function $g$ is orthogonal to the constants in $L^2(\pi)$
when the chain is reversible, and in general $\sum_x\pi_xg(x)=\sum_x\pi_x(Pg)(x)=\lambda\sum_x\pi_xg(x)$ forces
$\sum_x\pi_xg(x)=0$. Hence $\E_ig(X_U)=0$ and
\[
1=(1-\lambda)\,\E_i\sum_{s<U}g(X_s)\le(1-\lambda)\,\E_iU .
\]
For a reversible chain all eigenvalues are real; taking $\lambda=\lambda_2$ gives $\tmix\ge\min_U\E_iU\ge\tau_2$.
\end{proof}

\begin{proof}[Proof of the upper bound in Theorem~\ref{thm:main}]
Let $X$ have $m\le n$ states and let $\hX$ be an induced chain. We claim that
\begin{equation}\label{eq:chain}
\tmix(\hX)\ \le\ \tau_0(\hX)\ \le\ \tau_0(X)\ \le\ (m-1)\,\tau_2(X)\ \le\ (m-1)\,\tmix(X).
\end{equation}
For the first inequality, choose a state $J$ of $\hX$ with distribution $\hpi$, independently of the chain, and stop
at the hitting time of $J$; the chain is then in the stationary distribution, and the mean of this stopping time is
$\tau_0(\hX)$ from every starting state~\cite[Section~4.3]{AF}. The second inequality is part of the contraction
principle~\cite[Proposition~4.44]{AF}: by the random target lemma
$\tau_0=\tfrac12\sum_{x,y}\pi_x\pi_y(\E_xT_y+\E_yT_x)$, the commute time $\E_xT_y+\E_yT_x$ equals $w$ times the
effective resistance between $x$ and $y$, and gluing vertices does not increase effective resistances. The third is
the eigentime identity $\tau_0=\sum_{k\ge2}(1-\lambda_k)^{-1}$~\cite[Proposition~3.13 and Lemma~4.24]{AF}. The fourth
is Lemma~\ref{lem:relax}. If $\hX$ has one state then $\tmix(\hX)=0$, and $\tmix(X)\ge\tau_2(X)\ge\frac12$, so the
ratio is well defined; hence $K(n)\le n-1$.
\end{proof}

\begin{remark}\label{rem:pimin}
Write $\pi_*=\min_x\pi_x$. In discrete time $\tmix\le\frac{4e}{e-1}\,\tau_2\,(1+\frac12\log(1/\pi_*))$
by~\cite[Lemma~4.23]{AF}. Since $\tau_2(\hX)\le\tau_2(X)$ and $\min_a\hpi_a\ge\pi_*$, Lemma~\ref{lem:relax} gives
\[
\tmix(\hX)\ \le\ \frac{4e}{e-1}\Bigl(1+\frac12\log\frac1{\pi_*}\Bigr)\tmix(X),
\]
which is the form taken for $\tmix$ by the bound of Hermon and Peres~\cite[Remark~1.8]{HP}. Together with
Theorem~\ref{thm:main}, the factor is at most $\min\bigl(n-1,\ \frac{4e}{e-1}(1+\frac12\log(1/\pi_*))\bigr)$. For the
simple random walk on a graph with bounded degrees $\log(1/\pi_*)$ is of order $\log n$, and the examples
of~\cite{HP} show that this order is attained. In the comb of Section~\ref{sec:comb} with the parameters of
Table~\ref{tab:comb}, $\log(1/\pi_*)$ is of order $N\log t$.
\end{remark}

\section{The comb}\label{sec:comb}
We use the following elementary fact about hitting times across a bridge.

\begin{lemma}\label{lem:bridge}
Let $x\ne y$ be states with $w_{xy}>0$ such that every path from $x$ to $y$ in the network uses the edge $\{x,y\}$,
and let $S$ be the set of states that can be reached from $x$ without visiting $y$. Then
$\E_xT_y=w(S)/w_{xy}$, where $w(S)=\sum_{u\in S}w_u$.
\end{lemma}
\begin{proof}
Up to time $T_y$ the walk started at $x$ stays in $S$, so $\E_xT_y$ is the same for the network on $S\cup\{y\}$ in
which $y$ keeps only the edge to $x$. There the stationary probability of $y$ is $w_{xy}/(w(S)+w_{xy})$, so the mean
return time to $y$ is $1+w(S)/w_{xy}$, and it also equals $1+\E_xT_y$.
\end{proof}

Fix an integer $N\ge1$ and real numbers $r\ge2$, $\alpha>0$, $h\ge1$ and $\ell\ge0$. The \emph{comb} is the network
on the $2N+2$ states $z_0,\dots,z_N$, $y_0,\dots,y_{N-1}$, $c$ with the weights
\[
w_{z_kz_{k+1}}=r^k,\quad w_{y_kz_{k+1}}=\alpha r^k,\quad w_{y_ky_k}=(h-1)\alpha r^k\quad(0\le k<N),\qquad
w_{z_Nc}=r^N,\quad w_{cc}=\ell,
\]
and all other weights zero. The spine $z_0,\dots,z_N,c$ carries conductances that grow geometrically towards $c$; the
tooth $y_k$ hangs from $z_{k+1}$, and the walk leaves it with probability $1/h$ at each step. The map $f$ sends $z_k$
and $y_k$ to the same point $b_k$ for $0\le k<N$, sends $z_N$ to $b_N$ and $c$ to $c$. Put
\[
\delta=\frac{2+\alpha+\alpha h}{r-1},\qquad \rho=\frac{w_c}{r^N}=1+\frac{\ell}{r^N},\qquad
M=h+(N+1)(1+\delta),\qquad B=h+(N+2)(1+\delta).
\]

\begin{proposition}\label{prop:comb}
For the comb $X$ and its induced chain $\hX$,
\begin{enumerate}
\item[(i)] $\max_x\E_xT_c\le M$;
\item[(ii)] $\tmix(X)\le\max(B,\ M+B-\rho)$, so that $\tmix(X)\le B$ when $\rho\ge M$;
\item[(iii)] $\hE_{b_0}\widehat{T}_c\ge Nh\alpha/(1+\alpha)$ and $\tmix(\hX)\ge\pi_c\,Nh\alpha/(1+\alpha)$.
\end{enumerate}
\end{proposition}

\begin{proof}
The comb is a tree with loops. The vertex weights are $w_{z_0}=1$, $w_{z_k}=r^{k-1}(1+r+\alpha)$ for
$1\le k\le N$ and $w_{y_k}=h\alpha r^k$. Write $z_{N+1}=c$. For $0\le k\le N$ let $S_k$ be the total weight of the
states $z_0,\dots,z_k,y_0,\dots,y_{k-1}$, which are those separated from $c$ by the edge $\{z_k,z_{k+1}\}$. Since
$(1+r+\alpha+h\alpha)/(r-1)=1+\delta$,
\begin{equation}\label{eq:Sk}
S_k=1+(1+\delta)(r^k-1)\le(1+\delta)\,r^k .
\end{equation}

(i) By Lemma~\ref{lem:bridge}, $\E_{z_k}T_{z_{k+1}}=S_k/r^k\le1+\delta$ and $\E_{y_k}T_{z_{k+1}}=h$. Adding along the
path to $c$ gives $\E_{z_k}T_c\le(N+1)(1+\delta)$ and $\E_{y_k}T_c\le h+N(1+\delta)$.

(ii) Write $H(x,j)=\E_xT_j$ and $d(x,j)=H(x,j)-\E_\pi T_j=\sum_k\pi_k\bigl(H(x,j)-H(k,j)\bigr)$; by~\eqref{eq:lw},
$\tmix(X)=\max_{x,j}d(x,j)$. For three vertices $x,k,j$ of a tree let $m(x,k,j)$ be the vertex common to the three
paths joining them. A walk from $x$ to $j$ passes through $m=m(x,k,j)$, so $H(x,j)=H(x,m)+H(m,j)$, and likewise from
$k$; hence
\begin{equation}\label{eq:median}
H(x,j)-H(k,j)=H(x,m)-H(k,m)\le H(x,m).
\end{equation}
If $j=c$ then $d(x,c)\le H(x,c)\le M$ by (i).

Let $x=c$ and $j\ne c$. Define the level of $z_p$ to be $p$ and the level of $y_p$ to be $p+1$, and let $q$ be the
level of $j$. By Lemma~\ref{lem:bridge}, $\E_{z_{t+1}}T_{z_t}=(w-S_t)/r^t\le w\,r^{-t}$, so for $0\le p\le N$
\begin{equation}\label{eq:down}
H(c,z_p)\le w\sum_{t\ge p}r^{-t}=\frac{r}{r-1}\,w\,r^{-p}\le 2w\,r^{-p},
\end{equation}
and $\E_{z_{p+1}}T_{y_p}=(w-w_{y_p})/(\alpha r^p)\le w/(\alpha r^p)$. By~\eqref{eq:median},
$d(c,j)\le\sum_{k\ne c}\pi_kH(c,m(c,k,j))$. If the level $p$ of $k$ exceeds $q$ then $m(c,k,j)=z_p$; the two states
of level $p$ have total weight $r^{p-1}(1+r+\alpha+h\alpha)$, so by the geometric series in~\eqref{eq:down} they
contribute at most $(1+r+\alpha+h\alpha)/(r-1)=1+\delta$, and there are $N-q$ such levels. If the level of $k$ is at
most $q$ and $k\ne j$ then $m(c,k,j)=z_q$; the same holds for $k=j$ when $j=z_q$, while for $k=j=y_{q-1}$ the median
is $j$ and $H(c,j)=H(c,z_q)+H(z_q,y_{q-1})$. The states of level at most $q$, with $j$ among them, have total weight
$S_q$, so by~\eqref{eq:Sk} and~\eqref{eq:down} the terms $\pi_kH(c,z_q)$ contribute at most $2(1+\delta)$. If
$j=y_{q-1}$ there is in addition the term $\pi_jH(z_q,y_{q-1})$, which is at most $w_{y_{q-1}}/(\alpha r^{q-1})=h$.
Altogether
\[
d(c,j)\le h+(N-q+2)(1+\delta)\le B .
\]

Let $x\ne c$ and $j\ne c$. Since $c$ is a leaf, $m=m(x,c,j)\ne c$, and $m$ lies on the path from $x$ to $c$ and on the
path from $c$ to $j$. By~\eqref{eq:median}, (i) and Lemma~\ref{lem:bridge},
\[
H(x,j)-H(c,j)=H(x,m)-H(c,m)\le H(x,c)-H(c,z_N)\le M-\rho ,
\]
so $d(x,j)=H(x,j)-H(c,j)+d(c,j)\le M-\rho+B$. As $d(c,c)\le0$ and $M\le B$, this proves (ii).

(iii) The induced chain is a random walk on the path $b_0,\dots,b_N,c$ with $\hw_{b_kb_{k+1}}=(1+\alpha)r^k$ for
$k<N$, $\hw_{b_Nc}=r^N$ and loops. By Lemma~\ref{lem:bridge}, for $k<N$
\[
\hE_{b_k}\widehat{T}_{b_{k+1}}\ \ge\ \frac{\hw_{b_k}}{(1+\alpha)r^k}\ \ge\ \frac{w_{y_k}}{(1+\alpha)r^k}
=\frac{h\alpha}{1+\alpha},
\]
and the sum over $k$ gives the first assertion. On a path $\hE_{b_0}\widehat{T}_c\ge\hE_a\widehat{T}_c$ for every
$a$, so by~\eqref{eq:lw}
\[
\tmix(\hX)\ \ge\ \sum_a\hpi_a\bigl(\hE_{b_0}\widehat{T}_c-\hE_a\widehat{T}_c\bigr)\ \ge\ \hpi_c\,\hE_{b_0}\widehat{T}_c ,
\]
and $\hpi_c=\pi_c$.
\end{proof}

For the parameters used below, where $r$ is much larger than $\alpha h$, a walk on the spine of the comb rarely
enters a tooth, and the teeth cost a walk about one sojourn, of mean $h$, on its way to $c$. After lumping, every
point of the spine below $z_N$ inherits the loop of a tooth, and when $\alpha$ is large the delay is about $h$ at
each of the $N$ levels.

\begin{proof}[Proof of the lower bound in Theorem~\ref{thm:main}]
The identity map gives $K(n)\ge1$. Let $N\ge1$ and $t\ge2$, and take $\alpha=h=t$, $r=1+t^3$ and $\ell$ so large
that $\rho\ge M$ and $\pi_c\ge1-1/t$. Then $\delta=(2+t+t^2)/t^3\le1$ and, by Proposition~\ref{prop:comb},
\[
\frac{\tmix(\hX)}{\tmix(X)}\ \ge\ \Bigl(1-\frac1t\Bigr)\frac{N\,t^2}{(1+t)\bigl(t+(N+2)(1+\delta)\bigr)},
\]
which tends to $N$ as $t\to\infty$. The comb has $2N+2$ states, so $K(n)\ge N$ whenever $n\ge2N+2$.
\end{proof}

Table~\ref{tab:comb} gives the values for $\alpha=h=t=256$, $r=1+t^3$ and $w_c=(N+2)t^2(w-w_c)$, computed in
rational arithmetic from~\eqref{eq:lw} and rounded.

\begin{table}[h]
\centering
\begin{tabular}{rrrrr}
\toprule
$N$ & states & $\tmix(X)$ & $\tmix(\hX)$ & ratio\\
\midrule
1 & 4 & 257.004 & 256.012 & 0.9961\\
2 & 6 & 258.008 & 511.019 & 1.9806\\
4 & 10 & 260.016 & 1021.035 & 3.9268\\
8 & 18 & 264.031 & 2041.066 & 7.7304\\
\bottomrule
\end{tabular}
\caption{The comb with $\alpha=h=256$ and its induced chain.}\label{tab:comb}
\end{table}

\begin{remark}\label{rem:four}
The comb with $N=1$ has four states, and with the parameters of Table~\ref{tab:comb} its ratio is below one. A
different chain on four states has ratio above
one: take $I=\{0,1,2,3\}$ with $w_{02}=w_{12}=1$, $w_{13}=109$, $w_{23}=54$, $w_{11}=117$, $w_{33}=10000$, and lump
the states $2$ and $3$. Then
\[
\tmix(X)=\frac{149731115}{63193903}=2.3693\ldots,\qquad \tmix(\hX)=\frac{346900}{114917}=3.0187\ldots,
\]
both attained in~\eqref{eq:lw} at $(i,j)=(1,0)$, and the ratio is $419679620/329408453=1.2740\ldots$. Hence
$K(4)>1.27$.
\end{remark}

\section{Hitting times of sets}\label{sec:hit}
For a non-empty proper subset $Z$ of $I$ and $x,y\notin Z$ let $G_Z(x,y)$ be the expected number of visits to $y$ before $T_Z$
for the walk started at $x$, and let $R(y\lra Z)$ be the effective resistance between $y$ and $Z$ in the network with
conductances $w_{uv}$, $u\ne v$. Then~\cite[Section~2.2]{LP}
\begin{equation}\label{eq:green}
\E_xT_Z=\sum_{y\notin Z}G_Z(x,y),\qquad G_Z(x,y)=\Prob_x(T_y<T_Z)\,G_Z(y,y),\qquad G_Z(y,y)=w_y\,R(y\lra Z).
\end{equation}
The last identity holds in the presence of loops: the walk at $y$ reaches $Z$ before returning to $y$ with
probability $1/(w_yR(y\lra Z))$, and a loop at $y$ changes $w_y$ only.

\begin{proof}[Proof of Proposition~\ref{prop:hit}]
The network of $\hX$ is obtained from that of $X$ by gluing together the vertices in each set $f^{-1}(b)$; edges
inside such a set become loops. By Rayleigh's monotonicity principle~\cite[Section~2.4]{LP} gluing does not increase
effective resistances, so $\widehat{R}(b\lra\hA)\le R(y\lra A)$ for every $b\in\hI\setminus\hA$ and every
$y\in f^{-1}(b)$. With~\eqref{eq:green} applied to both chains,
\begin{align*}
\hE_a\widehat{T}_{\hA}\ &\le\ \sum_{b\notin\hA}\hw_b\,\widehat{R}(b\lra\hA)
=\sum_{b\notin\hA}\ \sum_{y\in f^{-1}(b)}w_y\,\widehat{R}(b\lra\hA)\\
&\le\ \sum_{y\notin A}w_y\,R(y\lra A)=\sum_{y\notin A}G_A(y,y)\ \le\ \sum_{y\notin A}\E_yT_A .
\end{align*}
For the second statement take the comb with $A=\{c\}$, $\hA=\{c\}$ and $a=b_0$. By
Proposition~\ref{prop:comb}(i) and~(iii),
\[
\frac{\hE_{b_0}\widehat{T}_c}{\max_y\E_yT_c}\ \ge\ \frac{Nh\alpha}{(1+\alpha)\bigl(h+(N+1)(1+\delta)\bigr)},
\]
and with $\alpha=h=t$, $r=1+t^3$ the right-hand side tends to $N$ as $t\to\infty$.
\end{proof}

The program in the repository named below evaluates both sides of every inequality of Proposition~\ref{prop:comb} in
rational arithmetic for $N\le8$ and $\alpha=h=t\le256$, and does the same for~\eqref{eq:chain} with $\tau_2$ omitted
and for Proposition~\ref{prop:hit} on $400$ random networks with three to six states and random partitions.

\section*{Acknowledgements and funding}
This work was done by multiple AI agents, Claude (Anthropic) and Codex (OpenAI), connected to Math
Brain~\cite{MathBrainRepo} and working together inside the Brain framework~\cite{BrainFramework}, a shared, persistent
memory in which the agents record claims, proofs, failed attempts and methods, retrieve one another's results,
coordinate who works on what, and audit each other's work. For more on the details behind the Brain and Math Brain, see
the Math Brain paper~\cite{MathBrainPaper}, the Math Brain repository~\cite{MathBrainRepo} and the Brain framework
repository~\cite{BrainFramework}.

The research presented in this paper was supported by the European Research Council (ERC) under the European Union's
Horizon Europe research and innovation programme (grant agreement No.~101041711), by the Simons Foundation as part of
the Collaboration on the Mathematical and Scientific Foundations of Deep Learning, by Heights Labs, by Convex Nexus
Capital, by the Israel Science Foundation (grant number 2258/19), and by the Israel Science Foundation (ISF Grant
4101/25). Computational resources were provided in part by Google Cloud through its education credits program (credit
530550103).

\section*{Declaration of competing interest}
The author holds equity in Heights Labs and in Convex Nexus Capital and has a management role at Convex Nexus Capital.
These entities are included in the funding declaration above.

\section*{Declaration of generative AI and AI-assisted technologies}
The paper comes out of Math Brain, an ongoing project on the use of AI systems for research mathematics, and its
results were developed with the Brain~\cite{BrainFramework} and Math Brain~\cite{MathBrainRepo}. In October 2026 an
AI agent connected to the Brain, working from the project's list of open problems in probability, found examples in
which lumping increases $\tmix$, among them the example of Remark~\ref{rem:four}. In November 2026 a Claude Fable
agent (Anthropic) obtained Theorem~\ref{thm:main}, Lemma~\ref{lem:relax} and Propositions~\ref{prop:hit}
and~\ref{prop:comb} with their proofs, wrote and ran the programs and prepared the manuscript. A Codex agent (OpenAI) then read the proofs and recomputed
the examples with programs of its own. These checks are not peer review. The author directed the work and is responsible for the content of the paper and for decisions
concerning its dissemination.

\section*{Code and data availability}
Programs and outputs:\\
\url{https://github.com/yspennstate/lumping-mixing-time}.

The paper is part of the Math Brain project, whose public repository is
\url{https://github.com/yspennstate/math-brain-public}~\cite{MathBrainRepo}.


\begin{thebibliography}{9}
\bibitem{Aldous82} D.~J.~Aldous, Some inequalities for reversible Markov chains, \emph{J. London Math. Soc.} (2)
\textbf{25} (1982), 564--576. \href{https://doi.org/10.1112/jlms/s2-25.3.564}{doi:10.1112/jlms/s2-25.3.564}.
\bibitem{AF} D.~Aldous and J.~A.~Fill, \emph{Reversible Markov Chains and Random Walks on Graphs}, unfinished
monograph, 2002, recompiled 2014, \url{https://www.stat.berkeley.edu/~aldous/RWG/book.html}.
\bibitem{HP} J.~Hermon and Y.~Peres, On sensitivity of mixing times and cutoff, \emph{Electron. J. Probab.}
\textbf{23} (2018), paper no.~25, 1--34. \href{https://doi.org/10.1214/18-EJP154}{doi:10.1214/18-EJP154};
arXiv:1610.04357.
\bibitem{LW95} L.~Lov\'asz and P.~Winkler, Efficient stopping rules for Markov chains, in \emph{Proceedings of the
Twenty-Seventh Annual ACM Symposium on Theory of Computing (STOC '95)}, ACM, New York, 1995, 76--82.
\href{https://doi.org/10.1145/225058.225086}{doi:10.1145/225058.225086}.
\bibitem{LW98} L.~Lov\'asz and P.~Winkler, Mixing times, in \emph{Microsurveys in Discrete Probability}, DIMACS Ser.
Discrete Math. Theoret. Comput. Sci. \textbf{41}, Amer. Math. Soc., Providence, RI, 1998, 85--133.
\href{https://doi.org/10.1090/dimacs/041/06}{doi:10.1090/dimacs/041/06}.
\bibitem{LP} R.~Lyons and Y.~Peres, \emph{Probability on Trees and Networks}, Cambridge Series in Statistical and
Probabilistic Mathematics \textbf{42}, Cambridge University Press, New York, 2016.
\href{https://doi.org/10.1017/9781316672815}{doi:10.1017/9781316672815}.

\bibitem{MathBrainRepo} Y.~Shmalo, \emph{Math Brain: public snapshot, code and tests}, 2026,
\url{https://github.com/yspennstate/math-brain-public}.
\bibitem{MathBrainPaper} Y.~Shmalo, \emph{Math Brain: shared, auditable memory for multi-agent mathematical research},
October 2026, \url{https://github.com/yspennstate/math-brain-public/tree/main/papers/math-brain-2026-10-24}.
\bibitem{BrainFramework} Y.~Shmalo, \emph{Brain Framework}, public software repository, 2026,
\url{https://github.com/yspennstate/brain-framework-public}.

\end{thebibliography}
\end{document}
